\documentclass[11pt]{article}

\usepackage[margin=1in]{geometry}
\usepackage{times}
\usepackage{amsmath}
\usepackage{booktabs}
\usepackage{graphicx}
\usepackage{hyperref}
\usepackage{url}
\usepackage{natbib}
\usepackage{xcolor}

\hypersetup{colorlinks=true, linkcolor=blue!50!black, citecolor=blue!50!black,
            urlcolor=blue!50!black}

\title{Ranking Without Resolution:\\
Most LLM Guardrail Benchmark Comparisons\\
Are Not Statistically Significant}

\author{
  Saurabh Kumbhar\\
  \texttt{fluiqai@gmail.com}\\
  \small Independent
}

\date{\today}

\begin{document}
\maketitle

\begin{abstract}
Vendor-published comparisons of LLM guardrails routinely report ranked tables of
recall, precision and F1 across a few dozen to a few hundred test cases. We show
that at these sample sizes the rankings are largely unresolvable. We benchmark
eight guardrails, six independent products, one thirty-line regular-expression
control and one belonging to the author, on 949 cases across four corpora, and
subject every pairwise comparison to McNemar's exact test on paired per-case
outcomes. On a 49-case adversarial suite of the kind vendors commonly publish,
only 11 of 28 pairs separate at $p < 0.05$; the four highest-scoring systems are
mutually indistinguishable, and the leading system is not distinguishable from
the regular-expression control ($p = 0.065$). On a 300-case jailbreak corpus, a
commercial detector achieving 97.3\% recall is not distinguishable from that same
control ($p = 0.248$), because an 83.3\% false-alarm rate cancels the advantage in
overall correctness. Monte Carlo power analysis over the observed disagreement
patterns indicates that separating two competent guardrails typically requires
between 1,600 and more than 6,400 cases, one to two orders of magnitude beyond
common practice. We argue that reporting a rank order without an interval or a
paired test is the central methodological defect in this literature, that a naive
control is the cheapest available diagnostic, and that the finding applies to the
author's own second-place results as forcefully as to anyone else's. Harness,
corpora, adapters and analysis code are released under MIT.
\end{abstract}

\section{Introduction}

An LLM guardrail is a classifier placed around a language model that decides
whether a given string should be blocked. Guardrails are sold, adopted and
compared, and the comparisons are usually published by one of the parties being
compared. The genre has a recognisable shape: a table of products, a column of
recall or F1 scores, a bolded row belonging to the publisher.

The methodological complaint usually levelled at these documents is bias, and it
is a fair complaint. Corpus selection, threshold choice, which configuration of
one's own product to enter and which competitor settings count as ``stock'' are
all degrees of freedom, and they reliably point one direction. But bias is
difficult to prove and easy to deny, and a publisher can respond to it with
disclosure and good intentions.

We argue that a second and more tractable defect is more damaging: the sample
sizes in common use cannot resolve the differences being reported, and almost no
published comparison says so. A table with eight rows and three decimal columns
communicates an ordering. If the ordering is not statistically distinguishable
from a permutation of itself, the table is not a weak result. It is a
misleading one, regardless of who published it and regardless of their care.

This paper makes three contributions.

\begin{enumerate}
\item We report a guardrail benchmark of eight systems on 949 cases across four
corpora, with Wilson confidence intervals on every rate and McNemar's exact test
on every pair. We find that a majority of pairwise comparisons on the smallest
corpus, and a meaningful minority on the larger ones, are not resolvable.

\item We show that the two most striking headline numbers in our own results are
artefacts of reporting recall without a paired test. A commercial detector with
the highest jailbreak recall in the study is statistically indistinguishable from
a thirty-line regular expression that detects almost nothing.

\item We estimate, by Monte Carlo over observed disagreement patterns, the corpus
size actually required to resolve these comparisons, and find it to be one to two
orders of magnitude larger than current practice.
\end{enumerate}

The benchmark is published by a party with a product in it. That conflict is
real, is stated throughout, and is partially mitigated by the direction of the
findings: the central result undercuts the author's own second-place finishes as
much as any competitor's first-place one.

\section{Related work}

\paragraph{Statistical testing of classifier comparisons.}
The problem of deciding whether one classifier is better than another on a shared
test set is old and well understood. \citet{dietterich1998} evaluates five
statistical tests for comparing supervised learning algorithms and recommends
McNemar's test \citep{mcnemar1947} for the case where algorithms are compared on
a single fixed test set, precisely the design used by guardrail benchmarks.
\citet{demsar2006} extends the treatment to multiple classifiers over multiple
datasets. \citet{wilson1927} gives the score interval for a binomial proportion
that we use throughout, which is well behaved at the small samples and extreme
proportions where the normal approximation fails.

\paragraph{Statistical power in NLP evaluation.}
\citet{card2020} document that a large fraction of published NLP experiments are
underpowered, such that reported improvements are unlikely to be detectable even
if real, and argue for power analysis as routine practice. \citet{bowman2021}
survey the broader failure modes of NLU benchmarking, including saturation,
annotation artefacts and the gap between leaderboard position and capability. Our
finding is a specific instance of the concern raised by both: the guardrail
literature has adopted the leaderboard format without the accompanying statistics.

\paragraph{Guardrails and prompt injection.}
Indirect prompt injection against LLM-integrated applications was characterised
by \citet{greshake2023}, following earlier demonstrations of instruction override
\citep{perez2022}. \citet{liu2024} propose a framework and benchmark for prompt
injection attacks and defences. Practical guardrail systems in wide use include
Microsoft Presidio for PII detection, NVIDIA NeMo Guardrails for programmable
rails, and LLM Guard for input and output scanning, alongside commercial
offerings. Existing benchmarks in this space overwhelmingly evaluate the
\emph{input} side, asking whether an attack prompt is detected. We evaluate the
output side, which we describe next.

\section{Benchmark design}

\subsection{Task}

Existing guardrail evaluation asks whether a model can be induced to misbehave.
We ask a different and narrower question: given a string that a model is about to
emit, or that a user has just sent, would the guardrail stop it? No model is run
during evaluation. Every system receives the identical string and returns a
binary verdict.

The output-side framing is motivated by an observation that shaped the corpus.
Contemporary frontier models refuse most direct attacks. In preliminary testing
against \texttt{claude-haiku-4-5}, seven attack framings including every
credential-exfiltration attempt were refused. What escapes is duller: data the
model was legitimately given, and identifiers that survive inside a refusal. One
corpus case is a verbatim model refusal that names two customers by email address
while explaining that it will not name them. A benchmark that only tests whether
the attack prompt is caught cannot see this failure at all.

\subsection{Corpora}

Four corpora, 949 cases (Table~\ref{tab:corpora}). Two are written for this
benchmark; two are sampled from public datasets that neither we nor any vendor
curated. Approximately one third of every corpus is benign text constructed to
bait false positives: order numbers shaped like card numbers, git SHAs shaped
like secrets, and sentences that merely \emph{describe} an injection attempt.

\begin{table}[t]
\centering
\begin{tabular}{lrrl}
\toprule
Corpus & Cases & Block/Allow & Source \\
\midrule
Output leakage, adversarial & 49 & 32/17 & written for this benchmark \\
Output leakage, public PII & 300 & 200/100 & \texttt{ai4privacy/pii-masking-200k} \\
Prompt injection & 300 & 150/150 & \texttt{deepset/prompt-injections} \\
Jailbreak & 300 & 150/150 & \texttt{jackhhao/jailbreak-classification} \\
\bottomrule
\end{tabular}
\caption{Corpora. The adversarial suite is deliberately of the size typical of
published vendor comparisons, which is part of what is under study.}
\label{tab:corpora}
\end{table}

The adversarial suite covers leak shapes a PII dataset does not contain:
obfuscated identifiers, credential formats, injected instructions echoed into
output, and refusals that leak while refusing. All of its values are synthetic.

For the public PII corpus, positives are rows containing at least one strong
identifier, and negatives are the \emph{same rows} with identifiers replaced by
placeholders, so both halves share a topic and register distribution and a
detector cannot score well by keying on subject matter.

\subsection{Systems evaluated}

Eight systems: three open source (LLM Guard, Presidio, NeMo Guardrails), three
commercial (Lakera Guard, Nightfall, AWS Comprehend), one belonging to the author
(Fluiq), and one control. Competitors run at stock settings with no tuning to
these cases.

The control is a thirty-line regular-expression matcher written in an afternoon.
Its purpose is diagnostic. A benchmark without a floor cannot distinguish
``the winner is strong'' from ``the field is weak,'' and if a naive baseline is
competitive with a commercial product, that is the most actionable fact
available to a reader.

Systems requiring credentials or an unavailable library are skipped and reported,
never scored zero. The harness refuses to rank any system that errors on more
than 2\% of a corpus: an early Nightfall run failed on 190 of 300 cases through
rate limiting, and scoring a rate limit as a detection failure is a way to
publish a falsehood accidentally. That run was discarded, backoff was added, and
the system re-run.

\subsection{Scoring}

Scoring was fixed before any system ran. A case labelled \emph{block} that is
blocked is a catch; blocked \emph{allow} is a false alarm; allowed \emph{block}
is a miss. We report recall and false-alarm rate side by side, never recall
alone, since a system that blocks everything achieves perfect recall. For the
paired tests we reduce each case to a single boolean: the system did the right
thing, meaning it blocked what should be blocked and allowed what should be
allowed.

\section{Statistical method}

\subsection{Intervals}

We report 95\% Wilson score intervals \citep{wilson1927} for recall and
false-alarm rate. At $n = 49$ and proportions near $0.9$ the normal approximation
produces upper bounds exceeding 100\%, which is not a confidence interval.

\subsection{Paired testing}

Every system is handed the identical case list, so outcomes are paired. We use
McNemar's exact test \citep{mcnemar1947}, following \citet{dietterich1998}, which
conditions on the discordant pairs: let $b$ be the number of cases the first
system handles correctly and the second does not, and $c$ the reverse. Under the
null hypothesis that the two are equally accurate, $b \sim \text{Binomial}(b+c,
0.5)$, and we compute the two-sided exact $p$-value

\begin{equation}
p = 2 \sum_{i=0}^{\min(b,c)} \binom{b+c}{i} 2^{-(b+c)},
\end{equation}

capped at 1. We use the exact form rather than the $\chi^2$ approximation because
discordant totals here are frequently below 25.

The choice of a paired test is not incidental. A two-proportion $z$-test on
independent samples, which is what an implicit comparison of two accuracy numbers
amounts to, discards the pairing. The pairing is the most informative feature of
the design: what matters is not how many cases each system got right, but the
cases on which the two disagree.

\subsection{Power}

To estimate the corpus size required to resolve a comparison, we treat the
observed joint outcome distribution for a pair as ground truth, resample cases
with replacement at a range of corpus sizes, re-run the test, and record the
fraction of simulations reaching $p < 0.05$. We report the smallest grid point
reaching 80\% power, over 2{,}000 simulations per point.

These cell probabilities are themselves estimated from small samples and carry
substantial uncertainty, particularly at $n = 49$. The figures should be read as
order-of-magnitude guidance for corpus design rather than precise requirements.
The direction of the result is robust to that uncertainty even where the exact
figure is not.

\section{Results}

\subsection{Adversarial suite, \texorpdfstring{$n = 49$}{n = 49}}

\begin{table}[t]
\centering
\small
\begin{tabular}{llrlrlr}
\toprule
System & Category & Recall & 95\% CI & FA & 95\% CI & F1 \\
\midrule
\texttt{llm-guard} & open source & 87.5 & [71.9, 95.0] & 17.6 & [6.2, 41.0] & \textbf{88.9} \\
\texttt{fluiq} & author & 81.2 & [64.7, 91.1] & 11.8 & [3.3, 34.3] & 86.7 \\
\texttt{aws-comprehend} & commercial & 84.4 & [68.2, 93.1] & 23.5 & [9.6, 47.3] & 85.7 \\
\texttt{nightfall} & commercial & 71.9 & [54.6, 84.4] & 5.9 & [1.0, 27.0] & 82.1 \\
\texttt{regex-baseline} & control & 62.5 & [45.3, 77.1] & 11.8 & [3.3, 34.3] & 74.1 \\
\texttt{lakera-guard} & commercial & 46.9 & [30.9, 63.6] & 5.9 & [1.0, 27.0] & 62.5 \\
\texttt{presidio} & open source & 43.8 & [28.2, 60.7] & 5.9 & [1.0, 27.0] & 59.6 \\
\texttt{nemo-guardrails} & open source & 43.8 & [28.2, 60.7] & 17.6 & [6.2, 41.0] & 57.1 \\
\bottomrule
\end{tabular}
\caption{Adversarial suite. All values are percentages. The intervals overlap
almost completely across the top five rows.}
\label{tab:adversarial}
\end{table}

Table~\ref{tab:adversarial} has the shape of a vendor comparison, and read as
one it supports a clear ordering. McNemar's test separates only 11 of the 28
pairs at $p < 0.05$. The four highest-scoring systems are mutually
indistinguishable: \texttt{llm-guard} versus \texttt{fluiq} gives $p = 1.000$
($b = 4$, $c = 3$), \texttt{llm-guard} versus \texttt{aws-comprehend} gives
$p = 0.791$, and \texttt{fluiq} versus \texttt{aws-comprehend} gives $p = 1.000$.

More pointedly, the leading system is not distinguishable from the control
($p = 0.065$, $b = 9$, $c = 2$), and neither is any other member of the top four.
A reader who concluded from Table~\ref{tab:adversarial} that \texttt{llm-guard}
is the strongest available output guardrail would be drawing an inference the
data does not license. A reader who concluded that the top four all beat a
thirty-line regular expression would also be drawing one.

\subsection{Public PII, \texorpdfstring{$n = 300$}{n = 300}}

At $n = 300$ resolution improves substantially: 24 of 28 pairs separate.
\texttt{aws-comprehend} (92.5\% recall, [88.0, 95.4]) is distinguishable from
every other system, and the control is distinguishable from everything. This is
what an adequately powered guardrail comparison looks like.

Four pairs remain unresolved, including the author's own second place over
\texttt{llm-guard} (81.0\% versus 80.0\%, $p = 0.473$). We report that ordering
in our published materials and it does not survive testing.

\subsection{Prompt injection and jailbreak, \texorpdfstring{$n = 300$}{n = 300}}

Only four systems implement input-side detection; the remainder are out of scope
and are not ranked. On prompt injection, \texttt{lakera-guard} reaches 94.7\%
recall at 16.7\% false alarms and is genuinely and distinguishably the strongest
system tested ($p < 0.001$ against all others). Everything else, including the
author's system, sits at F1 $\leq$ 54.5. Roughly a third of that corpus is not in
English, and that is where most misses concentrate.

The jailbreak corpus produces the study's sharpest illustration of the central
argument. \texttt{lakera-guard} achieves 97.3\% recall, [93.3, 99.0], the highest
recall figure anywhere in this paper. It also produces an 83.3\% false-alarm rate,
[76.6, 88.4]. Against the control, which achieves 0.7\% recall and detects
essentially nothing, McNemar returns $p = 0.248$ ($b = 145$, $c = 125$). The two
systems are not distinguishable in overall correctness.

We stress that this is not an argument that the control is a good guardrail, nor
that Lakera is a bad one; the operating points are radically different and a
deployment that tolerates false alarms may reasonably prefer high recall. It is
an argument about reporting. A headline of ``97.3\% jailbreak detection'' is
literally true, is the kind of number that appears in vendor material, and
conveys almost nothing about whether the system is more useful than trivial.

\subsection{Required corpus size}

\begin{table}[t]
\centering
\begin{tabular}{llrr}
\toprule
Corpus & Pair & Observed $b/c$ & Cases for 80\% power \\
\midrule
Adversarial ($n{=}49$) & \texttt{llm-guard} vs \texttt{aws-comprehend} & 8/6 & 1{,}600 \\
Adversarial ($n{=}49$) & \texttt{aws-comprehend} vs control & 10/5 & 400 \\
Adversarial ($n{=}49$) & control vs \texttt{nemo-guardrails} & 12/5 & 200 \\
PII ($n{=}300$) & \texttt{fluiq} vs \texttt{llm-guard} & 18/13 & 3{,}200 \\
PII ($n{=}300$) & \texttt{presidio} vs \texttt{nightfall} & 22/20 & $>$6{,}400 \\
Injection ($n{=}300$) & \texttt{llm-guard} vs \texttt{fluiq} & 20/17 & $>$6{,}400 \\
Jailbreak ($n{=}300$) & \texttt{lakera-guard} vs control & 145/125 & 3{,}200 \\
\bottomrule
\end{tabular}
\caption{Monte Carlo power estimates for unresolved pairs, 2{,}000 simulations
per point.}
\label{tab:power}
\end{table}

Table~\ref{tab:power} gives the estimated corpus size needed to resolve
representative unresolved pairs. The comparisons that matter most to a buyer,
those between two systems that are both broadly competent, are the most expensive
to resolve: separating the author's system from \texttt{llm-guard} on PII would
require roughly 3{,}200 cases, and several pairs exceed 6{,}400.

The pattern is a straightforward consequence of the arithmetic and is worth
stating plainly. Distinguishing a good guardrail from a bad one is cheap.
Distinguishing a good guardrail from another good guardrail is expensive, because
the two agree on almost every case and the test sees only the disagreements.
Vendor benchmarks are, almost by construction, comparisons among competent
systems, which is exactly the regime where a few hundred cases buy very little.

\section{Discussion}

\paragraph{The reporting defect is cheap to fix.}
Nothing in this paper requires new data collection. Confidence intervals and a
paired test over per-case outcomes are a few dozen lines of code, computable from
results any benchmark already produces. That they are absent from this literature
is a convention rather than a constraint.

\paragraph{A control is the highest-value single addition.}
Of everything reported here, the finding most likely to change a reader's
behaviour is that a thirty-line regular expression is statistically
indistinguishable from four commercial and open-source products on one corpus,
and from the highest-recall commercial detector on another. That diagnostic cost
an afternoon. We would suggest that a guardrail comparison without a naive
baseline should be treated as uninterpretable, since without a floor there is no
way to tell a strong winner from a weak field.

\paragraph{Recall without a false-alarm rate is not a result.}
This is well known and continues to be violated. We would go further: recall and
false-alarm rate together are still insufficient without a paired test, because
two systems at different operating points can produce very different-looking
pairs of numbers while being statistically identical in overall correctness, as
the jailbreak result demonstrates.

\paragraph{On publishing a benchmark one does not win.}
The author's system places second on both output-side corpora and third on both
input-side ones, and the central finding of this paper is that its second places
are not statistically meaningful. We regard the incentive structure here as worth
naming: a vendor benchmark that the vendor wins is close to uninformative,
because the space of defensible methodological choices is large enough that a
determined publisher can usually arrive there. The corresponding check available
to a reader is to ask whether the publisher reports a specific number that is
unflattering to them. Absence of such a number is itself the result.

\section{Limitations}

\paragraph{Conflict of interest.} This benchmark is published by a party with a
system in it. No methodology removes that. Mitigations: corpora, harness,
adapters and analysis code are public under MIT; scoring was fixed before any
system ran; competitors run at stock settings; every miss and false alarm is
listed by case identifier; a naive control is included; and the paper's principal
finding is adverse to the author's own results.

\paragraph{Single annotator.} Every label was assigned by one person, who also
works on one of the systems evaluated. There is no second annotator and therefore
no inter-annotator agreement statistic. This is the largest methodological gap in
the work. Per-case notes record the reasoning for contestable labels so that
disagreement can be specific, but that is a materially weaker guarantee than two
independent annotators and a reported $\kappa$. Some labels are explicitly
judgement calls: a company support inbox address is scored \emph{allow}, since
blocking it makes an assistant unusable, while ``I can confirm this person is a
customer, but I cannot share their SSN'' is scored \emph{block}, since it
confirms membership in a customer list. The jailbreak corpus counts a large
number of roleplay and persona prompts as benign, which is the single largest
source of label disagreement affecting these results.

\paragraph{Scale and coverage.} 949 cases is small, predominantly English, and
skewed toward PII and secrets. The 49-case adversarial suite is small enough that
its unresolved comparisons are unsurprising; we include it deliberately because
it is representative of published practice, but its size limits what can be
concluded from it beyond the resolution finding itself.

\paragraph{A string is not a trace.} Several systems, the author's included,
detect retrieval poisoning and tool-argument exfiltration from structured
execution context rather than from response text. Supplying only a string
understates them on the \texttt{rag\_poison\_echo} category. Correcting this
fairly requires a richer corpus format rather than an exception for one vendor.

\paragraph{Point-in-time measurement.} All systems were measured on 7 August
2026. Commercial detectors are opaque and change without notice, and several
lack a version identifier that a third party can pin. Any number here should be
read with its date attached.

\paragraph{Scope differences.} Presidio detects PII and does not claim to detect
API keys, so it scores zero on the secrets category. That is a scope difference
rather than a defect, which is why per-category results are reported and why no
single aggregate should be read alone.

\section{Conclusion}

Guardrail comparisons are published as ranked tables, and ranked tables assert an
ordering. We tested the orderings in a benchmark of eight systems and found that
a majority of the comparisons on a corpus of typical size do not survive a paired
significance test, that a thirty-line regular-expression control is
indistinguishable from several commercial products, and that resolving the
comparisons buyers actually care about would require one to two orders of
magnitude more data than is customary. The remedies are inexpensive: report
intervals, run a paired test, include a naive baseline, and state the date. We
have applied all four to our own results, where they eliminate the finding most
favourable to us.

\section*{Availability}

Harness, both hand-written corpora, public-dataset samplers, all adapters and the
analysis code in this paper are available under the MIT licence at
\url{https://github.com/SaurabhKumbhar24/guardrail-bench}. The public PII corpus
is not redistributed, as its upstream dataset card states no licence; the
repository ships a seeded sampler and a SHA-256 manifest of all 300 rows, and a
verification script that proves a rebuild is byte-identical to the corpus these
results were computed on.

\begin{thebibliography}{99}

\bibitem[Bowman and Dahl(2021)]{bowman2021}
Samuel R. Bowman and George E. Dahl.
\newblock What will it take to fix benchmarking in natural language understanding?
\newblock In \emph{Proceedings of NAACL-HLT}, 2021.

\bibitem[Card et al.(2020)]{card2020}
Dallas Card, Peter Henderson, Urvashi Khandelwal, Robin Jia, Kyle Mahowald, and
Dan Jurafsky.
\newblock With little power comes great responsibility.
\newblock In \emph{Proceedings of EMNLP}, 2020.

\bibitem[Dem\v{s}ar(2006)]{demsar2006}
Janez Dem\v{s}ar.
\newblock Statistical comparisons of classifiers over multiple data sets.
\newblock \emph{Journal of Machine Learning Research}, 7:1--30, 2006.

\bibitem[Dietterich(1998)]{dietterich1998}
Thomas G. Dietterich.
\newblock Approximate statistical tests for comparing supervised classification
learning algorithms.
\newblock \emph{Neural Computation}, 10(7):1895--1923, 1998.

\bibitem[Greshake et al.(2023)]{greshake2023}
Kai Greshake, Sahar Abdelnabi, Shailesh Mishra, Christoph Endres, Thorsten Holz,
and Mario Fritz.
\newblock Not what you've signed up for: Compromising real-world
LLM-integrated applications with indirect prompt injection.
\newblock In \emph{Proceedings of the 16th ACM Workshop on Artificial
Intelligence and Security (AISec)}, 2023.

\bibitem[Liu et al.(2024)]{liu2024}
Yupei Liu, Yuqi Jia, Runpeng Geng, Jinyuan Jia, and Neil Zhenqiang Gong.
\newblock Formalizing and benchmarking prompt injection attacks and defenses.
\newblock In \emph{Proceedings of USENIX Security}, 2024.

\bibitem[McNemar(1947)]{mcnemar1947}
Quinn McNemar.
\newblock Note on the sampling error of the difference between correlated
proportions or percentages.
\newblock \emph{Psychometrika}, 12(2):153--157, 1947.

\bibitem[Perez and Ribeiro(2022)]{perez2022}
F\'{a}bio Perez and Ian Ribeiro.
\newblock Ignore previous prompt: Attack techniques for language models.
\newblock In \emph{NeurIPS ML Safety Workshop}, 2022.

\bibitem[Wilson(1927)]{wilson1927}
Edwin B. Wilson.
\newblock Probable inference, the law of succession, and statistical inference.
\newblock \emph{Journal of the American Statistical Association},
22(158):209--212, 1927.

\end{thebibliography}

\end{document}
